%% academic-report-vi.tex
%% Template báo cáo học thuật tiếng Việt đầy đủ
%% Engine: XeLaTeX (BẮT BUỘC cho tiếng Việt)
%% Compile: xelatex → biber → xelatex → xelatex
%%
%% Cách compile:
%%   xelatex academic-report-vi.tex
%%   biber academic-report-vi
%%   xelatex academic-report-vi.tex
%%   xelatex academic-report-vi.tex

\documentclass[12pt, a4paper]{report}

%% ================================================================
%% PACKAGES
%% ================================================================

%% Encoding & Ngôn ngữ (XeLaTeX)
\usepackage{fontspec}
\usepackage{polyglossia}
\setmainlanguage{vietnamese}
\setotherlanguage{english}
\setmainfont{Times New Roman}
% Thay bằng font mã nguồn mở nếu không có Times New Roman:
% \setmainfont{FreeSerif}

%% Layout
\usepackage{geometry}
\geometry{
  a4paper,
  top=3cm, bottom=2.5cm,
  left=3.5cm, right=2cm
}

%% Header & Footer
\usepackage{fancyhdr}
\pagestyle{fancy}
\fancyhf{}
\fancyhead[L]{\small\leftmark}
\fancyhead[R]{\small\thepage}
\renewcommand{\headrulewidth}{0.4pt}

%% Hình ảnh
\usepackage{graphicx}
\usepackage{float}
\graphicspath{{figures/}}

%% Bảng
\usepackage{booktabs}
\usepackage{longtable}
\usepackage{tabularx}
\usepackage{multirow}

%% Math
\usepackage{amsmath}
\usepackage{amssymb}

%% Hyperlinks
\usepackage{hyperref}
\hypersetup{
  colorlinks=true,
  linkcolor=black,
  citecolor=blue,
  urlcolor=blue,
  pdftitle={Ứng dụng Học Sâu trong Phân Loại Văn Bản Tiếng Việt},
  pdfauthor={Nguyễn Văn An}
}

%% Bibliography
\usepackage[backend=biber, style=numeric, sorting=none]{biblatex}
\addbibresource{sample.bib}

%% Todo notes (xoá khi hoàn thiện)
\usepackage[colorinlistoftodos]{todonotes}

%% Spacing
\usepackage{setspace}
\onehalfspacing

%% ================================================================
%% BEGIN DOCUMENT
%% ================================================================
\begin{document}

%% ================================================================
%% 1. TITLE PAGE
%% ================================================================
\begin{titlepage}
  \centering

  {\Large\bfseries TRƯỜNG ĐẠI HỌC ABC\\[0.1cm]}
  {\large KHOA CÔNG NGHỆ THÔNG TIN\\[0.3cm]}
  \rule{\textwidth}{0.8pt}\\[1.5cm]

  {\Huge\bfseries ỨNG DỤNG HỌC SÂU\\[0.4cm]}
  {\Large\bfseries TRONG PHÂN LOẠI VĂN BẢN TIẾNG VIỆT\\[1cm]}

  {\large Báo cáo môn học: Xử lý Ngôn ngữ Tự nhiên\\[0.3cm]}
  {\large Mã học phần: IT4040\\[2cm]}

  \begin{tabular}{ll}
    \textbf{Sinh viên thực hiện:} & Nguyễn Văn An \\
    \textbf{MSSV:}                & 20215678 \\[0.3cm]
    \textbf{Giảng viên hướng dẫn:} & TS. Trần Thị Bình \\
    \textbf{Bộ môn:}              & Khoa học Máy tính \\
  \end{tabular}

  \vfill
  {\large Hà Nội, tháng 5 năm 2026}
\end{titlepage}

%% ================================================================
%% 2. ABSTRACT (TIẾNG VIỆT)
%% ================================================================
\chapter*{Tóm tắt}
\addcontentsline{toc}{chapter}{Tóm tắt}

Báo cáo này trình bày nghiên cứu về việc ứng dụng các kỹ thuật học sâu
(\textit{deep learning}) trong bài toán phân loại văn bản tiếng Việt.
Chúng tôi đề xuất một mô hình dựa trên kiến trúc Transformer được tinh chỉnh
(\textit{fine-tuned}) trên tập dữ liệu tiếng Việt gồm 50.000 mẫu văn bản
thuộc 10 chủ đề khác nhau.

Kết quả thực nghiệm cho thấy mô hình đề xuất đạt độ chính xác 92{,}7\%
trên tập kiểm tra, vượt trội hơn so với các phương pháp baseline truyền thống
như SVM (78{,}3\%) và LSTM (85{,}2\%). Đặc biệt, mô hình thể hiện hiệu suất
tốt với các văn bản có độ dài ngắn, điều mà các phương pháp trước đây
còn hạn chế \cite{Nguyen2024deep}.

\medskip
\noindent\textbf{Từ khóa:} học sâu, phân loại văn bản, tiếng Việt,
Transformer, xử lý ngôn ngữ tự nhiên.

%% Abstract tiếng Anh (tuỳ chọn)
\chapter*{Abstract}
\addcontentsline{toc}{chapter}{Abstract}

This report presents a study on applying deep learning techniques to Vietnamese
text classification. We propose a fine-tuned Transformer-based model trained on
a dataset of 50,000 Vietnamese text samples across 10 categories.

Experimental results demonstrate that our model achieves 92.7\% accuracy on the
test set, outperforming traditional baseline methods such as SVM (78.3\%) and
LSTM (85.2\%).

\medskip
\noindent\textbf{Keywords:} deep learning, text classification, Vietnamese,
Transformer, natural language processing.

%% ================================================================
%% 3. TABLE OF CONTENTS
%% ================================================================
\tableofcontents
\newpage

%% ================================================================
%% 4-5. LIST OF FIGURES & TABLES
%% ================================================================
\listoffigures
\addcontentsline{toc}{chapter}{Danh mục hình ảnh}

\listoftables
\addcontentsline{toc}{chapter}{Danh mục bảng biểu}
\newpage

%% ================================================================
%% 7. MAIN BODY
%% ================================================================

%% -------- CHƯƠNG 1: GIỚI THIỆU --------
\chapter{Giới thiệu}

\section{Đặt vấn đề}

Phân loại văn bản là một trong những bài toán nền tảng của xử lý ngôn ngữ
tự nhiên (NLP), với ứng dụng rộng rãi trong phân loại email spam, phân tích
cảm xúc, và tổ chức nội dung trực tuyến \cite{Goodfellow2016}.

Tuy nhiên, với tiếng Việt --- một ngôn ngữ đơn lập có hệ thống thanh điệu
phức tạp --- việc phân loại văn bản vẫn còn nhiều thách thức chưa được
giải quyết triệt để \cite{Le2024bert}. Các mô hình được huấn luyện cho
tiếng Anh thường không đạt hiệu suất tốt khi áp dụng trực tiếp cho tiếng Việt.

\section{Mục tiêu nghiên cứu}

Báo cáo này hướng đến ba mục tiêu chính:

\begin{enumerate}
  \item Xây dựng pipeline phân loại văn bản tiếng Việt sử dụng mô hình Transformer
  \item Đánh giá hiệu suất trên tập dữ liệu chuẩn
  \item So sánh với các phương pháp baseline (SVM, LSTM, CNN)
\end{enumerate}

\section{Phạm vi nghiên cứu}

Nghiên cứu tập trung vào:
\begin{itemize}
  \item Văn bản tin tức tiếng Việt (không xét văn bản mạng xã hội)
  \item 10 chủ đề: chính trị, kinh tế, thể thao, giải trí, công nghệ,
        sức khỏe, giáo dục, pháp luật, khoa học, môi trường
  \item Độ dài văn bản từ 50 đến 500 từ
\end{itemize}

\section{Cấu trúc báo cáo}

Phần còn lại của báo cáo được tổ chức như sau:
Chương~\ref{chap:theory} trình bày cơ sở lý thuyết.
Chương~\ref{chap:method} mô tả phương pháp đề xuất.
Chương~\ref{chap:results} trình bày và phân tích kết quả thực nghiệm.
Chương~\ref{chap:conclusion} kết luận và đề xuất hướng nghiên cứu tiếp theo.

%% -------- CHƯƠNG 2: CƠ SỞ LÝ THUYẾT --------
\chapter{Cơ sở lý thuyết}
\label{chap:theory}

\section{Học sâu và mạng nơ-ron}

Học sâu là một nhánh của học máy (\textit{machine learning}) sử dụng mạng
nơ-ron nhân tạo nhiều lớp để học các biểu diễn đặc trưng từ dữ liệu thô
\cite{LeCun2015deep}. Hàm kích hoạt ReLU được định nghĩa:

\begin{equation}
  \label{eq:relu}
  f(x) = \max(0, x)
\end{equation}

\section{Kiến trúc Transformer}

Transformer \cite{Vaswani2017attention} là kiến trúc mạng nơ-ron dựa trên
cơ chế self-attention, cho phép mô hình học được mối quan hệ giữa các từ
trong câu bất kể khoảng cách.

Cơ chế attention được tính như sau:

\begin{equation}
  \label{eq:attention}
  \text{Attention}(Q, K, V) = \text{softmax}\!\left(\frac{QK^T}{\sqrt{d_k}}\right) V
\end{equation}

trong đó $Q$, $K$, $V$ lần lượt là ma trận query, key và value;
$d_k$ là chiều của vector key.

\section{Các phương pháp baseline}

\begin{table}[htbp]
  \centering
  \caption{So sánh các phương pháp phân loại văn bản}
  \label{tab:methods}
  \begin{tabular}{lll}
    \toprule
    Phương pháp & Đặc điểm & Hạn chế \\
    \midrule
    SVM         & Hiệu quả với feature thủ công & Không học tự động \\
    LSTM        & Xử lý chuỗi tuần tự & Chậm, khó song song hóa \\
    CNN         & Nhanh, hiệu quả với n-gram & Bỏ qua thứ tự xa \\
    Transformer & Self-attention toàn cục & Tốn bộ nhớ $O(n^2)$ \\
    \bottomrule
  \end{tabular}
\end{table}

%% -------- CHƯƠNG 3: PHƯƠNG PHÁP --------
\chapter{Phương pháp nghiên cứu}
\label{chap:method}

\section{Tập dữ liệu}

Chúng tôi sử dụng tập dữ liệu VNews-50K gồm 50.000 bài báo tiếng Việt
thu thập từ các tờ báo điện tử lớn trong giai đoạn 2020--2024. Tập dữ liệu
được chia theo tỷ lệ 70/15/15 cho train/validation/test.

\begin{table}[htbp]
  \centering
  \caption{Thống kê tập dữ liệu VNews-50K}
  \label{tab:dataset}
  \begin{tabularx}{\textwidth}{Xrrr}
    \toprule
    Chủ đề & Train & Validation & Test \\
    \midrule
    Chính trị  & 3.500 & 750 & 750 \\
    Kinh tế    & 3.500 & 750 & 750 \\
    Thể thao   & 3.500 & 750 & 750 \\
    Giải trí   & 3.500 & 750 & 750 \\
    Công nghệ  & 3.500 & 750 & 750 \\
    \midrule
    Tổng       & 35.000 & 7.500 & 7.500 \\
    \bottomrule
  \end{tabularx}
\end{table}

\section{Mô hình đề xuất}

Chúng tôi fine-tune mô hình PhoBERT \cite{Le2024bert} trên tập dữ liệu
VNews-50K. Quá trình huấn luyện sử dụng các siêu tham số sau:

\begin{itemize}
  \item Learning rate: $2 \times 10^{-5}$
  \item Batch size: 32
  \item Số epoch: 10
  \item Optimizer: AdamW với weight decay $0.01$
\end{itemize}

%% Ví dụ thêm hình ảnh (bỏ comment khi có file ảnh)
%\begin{figure}[htbp]
%  \centering
%  \includegraphics[width=0.85\textwidth]{architecture.png}
%  \caption{Kiến trúc mô hình đề xuất}
%  \label{fig:architecture}
%\end{figure}

%% -------- CHƯƠNG 4: KẾT QUẢ --------
\chapter{Kết quả và thảo luận}
\label{chap:results}

\section{Kết quả thực nghiệm}

Bảng~\ref{tab:results} so sánh hiệu suất của các mô hình trên tập test.
Mô hình đề xuất đạt kết quả tốt nhất trên tất cả các chỉ số.

\begin{table}[htbp]
  \centering
  \caption{Kết quả so sánh hiệu suất các mô hình (tập test)}
  \label{tab:results}
  \begin{tabular}{lcccc}
    \toprule
    Mô hình      & Accuracy & Precision & Recall & F1-score \\
    \midrule
    SVM          & 78.3\%   & 0.781     & 0.783  & 0.779 \\
    CNN          & 83.1\%   & 0.829     & 0.831  & 0.828 \\
    LSTM         & 85.2\%   & 0.851     & 0.852  & 0.849 \\
    \textbf{Đề xuất (PhoBERT)} & \textbf{92.7\%} & \textbf{0.926} & \textbf{0.927} & \textbf{0.923} \\
    \bottomrule
  \end{tabular}
\end{table}

\section{Phân tích và thảo luận}

Kết quả cho thấy mô hình Transformer vượt trội đáng kể so với các phương
pháp truyền thống. Cải thiện lớn nhất được quan sát ở chủ đề ``Công nghệ''
(+9.2\% so với LSTM), nơi văn bản thường chứa nhiều thuật ngữ chuyên biệt.

Tuy nhiên, mô hình vẫn còn hạn chế với các văn bản rất ngắn (dưới 50 từ),
khi F1-score giảm xuống còn 85.4\%. \todo{Phân tích thêm trường hợp lỗi}

%% -------- CHƯƠNG 5: KẾT LUẬN --------
\chapter{Kết luận}
\label{chap:conclusion}

\section{Tóm tắt đóng góp}

Báo cáo này đã:
\begin{enumerate}
  \item Xây dựng thành công pipeline phân loại văn bản tiếng Việt
        dựa trên mô hình PhoBERT
  \item Đạt độ chính xác 92{,}7\% trên tập kiểm tra, vượt trội hơn
        các baseline tới 7{,}5 điểm phần trăm
  \item Phân tích điểm mạnh và hạn chế của mô hình đề xuất
\end{enumerate}

\section{Hướng nghiên cứu tiếp theo}

Trong tương lai, chúng tôi dự kiến:
\begin{itemize}
  \item Mở rộng tập dữ liệu lên 200.000 mẫu, bao gồm văn bản mạng xã hội
  \item Thử nghiệm với các mô hình ngôn ngữ lớn mới hơn
  \item Xây dựng ứng dụng web để phân loại văn bản thời gian thực
\end{itemize}

%% ================================================================
%% 8. BIBLIOGRAPHY
%% ================================================================
\printbibliography[title={Tài liệu tham khảo}]

%% ================================================================
%% 9. APPENDICES
%% ================================================================
\appendix

\chapter{Mã nguồn huấn luyện mô hình}
\label{app:code}

Đoạn code Python dưới đây minh họa quá trình fine-tune PhoBERT:

\begin{verbatim}
from transformers import AutoTokenizer, AutoModelForSequenceClassification
import torch

model_name = "vinai/phobert-base"
tokenizer = AutoTokenizer.from_pretrained(model_name)
model = AutoModelForSequenceClassification.from_pretrained(
    model_name, num_labels=10
)

# Fine-tuning loop
optimizer = torch.optim.AdamW(model.parameters(), lr=2e-5)
for epoch in range(10):
    for batch in train_dataloader:
        outputs = model(**batch)
        loss = outputs.loss
        loss.backward()
        optimizer.step()
        optimizer.zero_grad()
\end{verbatim}

\chapter{Danh sách tập dữ liệu}
\label{app:dataset}

Chi tiết nguồn dữ liệu được trình bày trong Bảng~\ref{tab:data-sources}.

\begin{table}[htbp]
  \centering
  \caption{Nguồn dữ liệu VNews-50K}
  \label{tab:data-sources}
  \begin{tabular}{lrl}
    \toprule
    Nguồn & Số bài & URL \\
    \midrule
    VnExpress  & 20.000 & \url{https://vnexpress.net} \\
    Tuổi Trẻ   & 15.000 & \url{https://tuoitre.vn} \\
    Thanh Niên & 15.000 & \url{https://thanhnien.vn} \\
    \midrule
    Tổng       & 50.000 & \\
    \bottomrule
  \end{tabular}
\end{table}

\end{document}
